\documentclass[10pt,a4paper]{amsart}
\usepackage[margin=19mm]{geometry}
\usepackage{amsmath,amssymb,mathtools}
\usepackage[scr=rsfs,cal=euler]{mathalfa}
\usepackage{xfrac}
\usepackage[unicode,hidelinks,bookmarks=false]{hyperref}
\newcommand{\ie}{i.e.\ }
\newcommand{\cf}{cf.\ }
\newcommand{\resp}{respectively\ }
\newcommand{\Subsection}{\subsection}
\providecommand{\nobreakdash}{\nobreak\hbox{-}}
\setlength{\emergencystretch}{2em}
\multlinegap0pt
\usepackage{bm}
\usepackage{bbm}
\usepackage{dsfont}
\usepackage{xspace}
\usepackage{cancel}
\usepackage{mathtools}
\usepackage{amsthm}
\makeatletter
\@ifundefined{thm}{\newtheorem{thm}{Thm}}{}
\@ifundefined{cor}{\newtheorem{cor}[thm]{Corollary}}{}
\@ifundefined{lem}{\newtheorem{lem}[thm]{Lemma}}{}
\makeatother
\newcommand{\N}{{\mathbb{N}}}
\newcommand{\R}{{\mathbb{R}}}
\newcommand{\cinf}{{\mathcal C}^{\infty}}
\newcommand{\map}{\operatorname{Map}}
\title[Derived Lie $\infty$-groupoids and algebroids in higher diff]{Derived Lie $\infty$-groupoids and algebroids in higher differential geometry}
\author{Qingyun Zeng}
\date{Source: arXiv:2210.05856v2 (CC BY 4.0)}
\begin{document}
\maketitle

\noindent\emph{Source and licence.} Derived Lie $\infty$-groupoids and algebroids in higher differential geometry. Version arXiv:2210.05856v2 (CC BY 4.0), distributed under the Creative Commons Attribution 4.0 International licence, as recorded in the arXiv licence metadata for this exact version and shown on the abstract page https://arxiv.org/abs/2210.05856v2. Full rights notice: licenses/arxiv-2210.05856-CC-BY-4.0.txt.\par\smallskip

\noindent\textit{Editorial scope.} Counted unit: the Lem whose source label is \texttt{lem1} in the article (source0.tex lines 6244--6245, source label \texttt{lem1}), delivered together with the article's own complete original proof of it (source0.tex lines 6247--6319, one \texttt{proof} environment of 73 lines). No mathematics is written, repaired or re-derived: statement and proof are the article's own text line for line. Required context retained uncounted: none. Every definition, display and label of those lines is retained unchanged. Omitted from this package and not redistributed: 1--6243, 6246--6246, 6320--7278. Nothing inside a retained excerpt is omitted or altered: every retained body is delivered exactly as the article's lines are written, so no omission receipt of any kind accompanies this package. Citations: the retained excerpts carry no citation command at all, so no bibliography entry of the article is transcribed and no reference is redistributed. Reference adaptation: none was needed and none was made; every reference inside the delivered unit resolves inside it, so no unresolved reference or citation is produced. Printed slips kept literal: no printed word, symbol, hypothesis or number of the counted statement or of its proof was altered. Layout and class: the article's own class is \texttt{amsart} with packages including mathpazo, avant, tensor, todonotes, biblatex, graphicx, amssymb, tikz-cd; the wrapper is the lane's pinned \texttt{amsart} editorial wrapper at 10pt on A4, which restates the theorem-like environments the retained text needs and adds the article's own macro definitions that the retained text needs (\texttt{\textbackslash N}, \texttt{\textbackslash R}, \texttt{\textbackslash cinf}, \texttt{\textbackslash map}), with extra wrapper packages (bm, bbm, dsfont, xspace, cancel, mathtools). The delivered numbering is the unit's own. Assets: no retained span contains a figure environment, a graphics-inclusion command, a PDF-page inclusion, a table or a TikZ picture; no image file is copied into this package, the package declares no asset dependency and no image file is redistributed with it. Licence: version arXiv:2210.05856v2 of this article is distributed under the Creative Commons Attribution 4.0 International licence, as recorded in the arXiv licence metadata for this exact version and shown on the abstract page https://arxiv.org/abs/2210.05856v2. A full rights notice is delivered at licenses/arxiv-2210.05856-CC-BY-4.0.txt. Characters: every retained excerpt is pure ASCII, so the delivered engine, XeLaTeX with its default Latin Modern text font, reports no missing character.
	 \begin{lem}\label{lem1} Let $h\in \cinf(\R)$ to be strictly positive for $x<0$ and 0 for $x \ge 0$, and $g\in \cinf(\R^2)$. Let $\cinf(\R^2)$ as a $\cinf(\R)$-module induced by the projection $\R^2 \to \R$. Suppose $hg =0$, then $g(x,y) = c(x)G(x, y)$, where $c(x) \in \cinf(\R)$ vanishes on $x\le 0$ and $G(x,y) \in \cinf(\R^2)$ .
	 	\end{lem}

 	\begin{proof}
 		Consider two sets $I = \{f\in \cinf(\R)|\text{$f=0$ when $x\le 0$ and $f> 0$ when $x> 0$} \}$, $J = \{f\in \map(\R^+, \R)| \text{$f(x)/x^n \to 0$ as $x\to 0$ for all $n>0$}\}$.
 		
 		\begin{lem} For any $f\in J$, there exists a $g\in I$ such that $f/g\to 0$ as $x\to 0$.
 			\end{lem}
 		\begin{proof}
 			Consider a bump function $r\in \cinf(\R)$ such that $r= 1$ for $x\le 0$ and $r=0$ for $x \ge 1$. Let $\{a_k\}$ be a monotonically decreasing sequence such that $\sum_i a_k < \infty$ and $a_k>0$ for all $k$. Define $\theta(x) = \sum_{i=1}^{\infty} r(x/a_i)$. Let 
 			\begin{equation}
 			g(x) = \begin{cases}
 			x^{\theta(x)} \quad &x>0\\
 			0 &x\le 0
 			\end{cases}.
 			\end{equation}
 			We claim that $g(x)$ satisfied the requirement in lemma. Outside any open neighborhood of $0$, there will be only finitely many non-zero summands in $\theta$, hence $g(x)$ is smooth and bounded outside any open neighborhood of $0$. Clearly $g(x)/x^n\to 0$ as $x\to 0$. We just need to check all derivative of $g$. Let $\theta_n(x) = \sum_{i=1}^n a_k$ and $g_n = x^{\theta_n}$, we have 
 			$$
 			\frac{d^l}{dx^l}g_n =\frac{d^l}{dx^l} x^{\theta_n} = \theta_n(x)\cdots(\theta_n(x)-l+1) x^{\theta_n(x) - l}
 			$$
 			for $x>0$. For each $n$, there exists an $\epsilon_n$ such that for $0<x<\epsilon_n$, $\frac{d^l}{dx^l} x^{\theta_{n+1}}<\frac{d^l}{dx^l} x^{\theta_n}$. For each $l$, if we pick $n$ large enough, e.g. $l >n$, then $\frac{d^l}{dx^l} x^{\theta_n} = 0$. Let $\epsilon>0$ be arbitrary, we want to show that there exists  $x_0>0$ and $N\in \N$ such that for all $n>N, 0<x<x_0$, $|\frac{d^l}{dx^l} x^{\theta_n(x)}-\frac{d^l}{dx^l} x^{\theta(x)}| < \epsilon$, that is $x^{\theta_n}\to x^{\theta}$ uniformly on $[0, x_0]$. Consider $n>l$ to be sufficiently large, then there exist $x_1> 0$ such that $\frac{d^l}{dx^l} x^{\theta_n} < \epsilon$ for $0\le x \le x_1$. Now
 			\begin{align*}
 				\frac{d^l}{dx^l} x^{\theta_{n+1}}& = \theta_{n+1}(x)\cdots(\theta_{n+1}(x)-l+1) x^{\theta_{n+1}(x) - l}\\
 				& =\frac{\theta_{n+1}(x)\cdots(\theta_{n+1}(x)-l+1)}{\theta_n(x)\cdots(\theta_n(x)-l+1)}\theta_n(x)\cdots(\theta_n(x)-l+1) x^{\theta_n(x) - l}x^{r(x/a_{n+1})} \\
 				& \le x^{\theta_{n+1}(x) - l}  \Bigg(\frac{\theta_{n+1}(x)-l+1}{\theta_{n}(x)-l+1} \Bigg)^l \frac{d^l}{dx^l} x^{\theta_{n}}   
 			\end{align*}
 		    Note that $\frac{\theta_{n+1}(x)-l+1}{\theta_{n}(x)-l+1}$ is monotonically decreasing as $n$ increases and as $x$ decreases. Let $\delta >0$ such that $(\frac{\theta_{n+1}(x)-l+1}{\theta_{n}(x)-l+1})^l < 1+\delta$ for all $0\le 0 \le x_1$, then we can find an $x_2<x_1$ such that $x^{r(x/a_{n+1})} <1/(1+\delta)$. Therefore, $\frac{d^l}{dx^l} x^{\theta_{n+1}} < \epsilon$ on $[0, x_2]$. By induction, we get  $\frac{d^l}{dx^l} x^{\theta_{k}} < \epsilon$ for all $k\ge n$ and $0\le x \le x_2$. 
 		    
 		    Now picking a sequence $\{x_m\}$ such that all $x_m \le x_2$ and $x_m \to 0$ monotonically, then $\lim_{m\to 0} g(x_m) = 0$ by continuity, and 
 		    $$
 		    \lim_{n\to \infty}\lim_{m\to 0}\frac{d^l}{dx^l}g_n(x_m) = \lim_{n\to \infty}\frac{d^l}{dx^l}g_n(0) = 0
 		    $$,
 		    then by the uniform convergence, we can change the order of limits and get
 			$$
 			\lim_{m\to 0}\lim_{n\to \infty}\frac{d^l}{dx^l}g_n(x_m) = \lim_{m\to 0}\frac{d^l}{dx^l}g(x_m) = 0
 			$$
 			
 			Now we have shown that $g\in J$. Since $f=0$ for all $x\le 0$, all derivatives of $f$ must vanish at infinite order at $0$, hence for each $n>0$, there exists a decreasing sequence $\{\epsilon_n\}$ and $\epsilon_k <1$ for all $k$ such that $|f|\le x^{n}$ for all $x\in [0, \epsilon_n]$. Now we just need to pick $a_k < \epsilon_{n+1}$ for all $k\ge n$ which ensure that $g(x)/x^{n-1}>1$ for $x \in (\epsilon_{n+1},\epsilon_n)$. Hence, $|f(x)/g(x)| \le x $  on $(\epsilon_{n+1},\epsilon_n)$ for all $n$, which implies $f/g \to 0$ as $x \to 0$.
 			
 	
% 			
% 			$$
% 			\frac{d^n}{dx^n}\log \big(g(x)\big) = \frac{d^n}{dx^n} \big(\theta(x)\log(x)\big)
% 			$$
% 			since $r(x)$ has compact support, $r^{(n)}(x) \le C_n x^{-n-1}$ for some $C_n$, we have
% 			
% 			$$
% 				\frac{d^n}{dx^n}\theta(x) = \sum_i^{\infty}2^{in}r^{(n)}(2^i x) \le C_n x^{-n-1}
% 			$$
% 			therefore, $\frac{d^n}{dx^n} \big(\theta(x)\log(x)\big) \sim O(x^{-n-2})$ for $n\ge 1$. On the other hand, $\frac{d^n}{dx^n}\log \big(g(x)\big)= \frac{d^{n-1}}{dx^{n-1}}(g'/g)$ consists of $g^{(n)}/g$ and a polynomial in $g^{(i)}/g$ for $i<n$ 
% 			$$
% 			\frac{d^{n-1}}{dx^{n-1}}\bigg(\frac{g'}{g}\bigg) =   \frac{1}{g} \left( g^{(n)} -\sum_{j=1}^{n-1} \binom{n-1}{j} \bigg(\frac{g'}{g}\bigg)^{(n-1-j)}g^{(j)} \right).
% 			$$ 
% 			by induction we see that $g^{(n)}/g \simeq O(x^{-n-2})
% 			
 			\end{proof}
 		\begin{cor}
 			Given a sequence $\{f_i\}$ in $J$, there exists a $g\in I$ such that $f_k/g \to 0$ as $x \to 0$ for all $k$.
 		\end{cor}
 	
 		\begin{cor}
 		Given a sequence $\{f_i\}$ in $J$, there exists a $g\in I$ such that $f_k/g^n \to 0$ as $x \to 0$ for all $k, n \in \N$.
 	\end{cor}
 
 Now consider $f_{ijk}(x) = \sup_{|y| \le k}\big(\frac{d^{i+j}}{dx^i dy^j}g(x,y)\big)$. Then by previous lemma, we can find an $a(x)$ such that $f_{ijk}/a^n \to 0$ as $x \to 0^+$ for all $i, j, k, n$. Now, take $$G(x, y)= \begin{cases}
                   g(x,y)/a(x)\quad & x >0\\
                   0 & x\le 0
            \end{cases}
 $$
 It suffices to verify $G^{(n)}(x,y)\to 0$ as $x\to 0^+$. Expand by quotient rules, we have
 $$
  			\bigg(\frac{g(x,y)}{a(x)}\bigg)^{(n)} =   \frac{1}{a} \left( g^{(n)} -\sum_{j=1}^{n} \binom{n}{j} \bigg(\frac{g}{a}\bigg)^{(n-1-j)}a^{(j)} \right).
 			$$
which goes to $0$ by induction. Hence, $g=aG$ is the desired factorization.                                                                                                                                                                                                                                                                                                                                                                                                                                                                                                                                                                                                          
 
\end{proof}

\end{document}
